\chapter{2024年上海卷（春）}

\section{填空题}

\begin{problem}
函数 \(y=\log_2 x\) 的定义域为\fillinblank{} .
\end{problem}

\begin{answer}
\(\left(0,+\infty\right)\)
\end{answer}

\begin{solution}
由对数真数必须大于零，对 \(y=\log_{2}x\) 得 \(x>0\)，故定义域为 \(\left(0,+\infty\right)\)，回代 \(x>0\) 时 \(\log_{2}x\) 有意义，而 \(x=0\) 处无定义，边界检查成立.
\end{solution}

\begin{problem}
直线 \(x-y+1=0\) 的倾斜角为\fillinblank{} .
\end{problem}

\begin{answer}
\(\frac{\pi}{4}\)
\end{answer}

\begin{solution}
将 \(x-y+1=0\) 化为 \(y=x+1\)，得斜率为 \(1\)，设倾斜角为 \(\theta\)，则 \(\theta\in\left[0,\pi\right)\) 且 \(\tan\theta=1\)，故 \(\theta=\frac{\pi}{4}\)，回代斜率与截距满足原方程.
\end{solution}

\begin{problem}
已知 \(\frac{z}{1+\i}=i\)，则 \(z\) 的共轭复数为\fillinblank{} .
\end{problem}

\begin{answer}
\(-1-\i\)
\end{answer}

\begin{solution}
由 \(\frac{z}{1+\i}=\i\) 得 \(z=\i\left(1+\i\right)=\i+\i^{2}=-1+\i\)，取共轭得 \(\overline{z}=-1-\i\)，回代 \(\frac{-1+\i}{1+\i}=\i\) 成立.
\end{solution}

\begin{problem}
\(\left(x-1\right)^6\) 展开式中 \(x^4\) 的系数为\fillinblank{} .
\end{problem}

\begin{answer}
\(15\)
\end{answer}

\begin{solution}
由二项式定理，\(\left(x-1\right)^{6}\) 的通项为 \(\mathrm C_{6}^{k}x^{k}\left(-1\right)^{6-k}\)，取 \(k=4\) 得含 \(x^{4}\) 的项为 \(\mathrm C_{6}^{4}x^{4}\left(-1\right)^{2}=15x^{4}\)，故 \(x^{4}\) 的系数为 \(15\).
\end{solution}

\begin{problem}
三角形 \(ABC\) 中，\(BC=2\)，\(A=\frac{\pi}{3}\)，\(B=\frac{\pi}{4}\)，则 \(AB=\)\fillinblank{} .
\end{problem}

\begin{answer}
\(\frac{3\sqrt{2}+\sqrt{6}}{3}\)
\end{answer}

\begin{solution}
由 \(A+B+C=\pi\)，\(A=\frac{\pi}{3}\)，\(B=\frac{\pi}{4}\) 得 \(C=\pi-\frac{\pi}{3}-\frac{\pi}{4}=\frac{5\pi}{12}\)，由正弦定理 \(\frac{AB}{\sin C}=\frac{BC}{\sin A}\)，\(BC=2\)，得 \(AB=\frac{2\sin\frac{5\pi}{12}}{\sin\frac{\pi}{3}}\)，而 \(\sin\frac{5\pi}{12}=\sin\left(\frac{\pi}{4}+\frac{\pi}{6}\right)=\frac{\sqrt{6}+\sqrt{2}}{4}\)，\(\sin\frac{\pi}{3}=\frac{\sqrt{3}}{2}\)，代入得 \(AB=\frac{3\sqrt{2}+\sqrt{6}}{3}\).
\end{solution}

\begin{problem}
已知 \(ab=1\)，\(4a^2+9b^2\) 的最小值为\fillinblank{} .
\end{problem}

\begin{answer}
\(12\)
\end{answer}

\begin{solution}
记 \(X=2a\)，\(Y=3b\)，则 \(4a^{2}+9b^{2}=X^{2}+Y^{2}\ge 2XY=12ab\)，已知 \(ab=1\)，故 \(4a^{2}+9b^{2}\ge 12\)，上式对任意实数 \(X\)，\(Y\) 成立，不需 \(a\)，\(b\) 同为正数，等号条件为 \(X=Y\) 即 \(2a=3b\)，与 \(ab=1\) 联立有实数解 \(a=\sqrt{\frac{3}{2}}\)，\(b=\sqrt{\frac{2}{3}}\)（同取相反数亦满足 \(2a=3b\) 与 \(ab=1\)），故等号可取，最小值为 \(12\).
\end{solution}

\begin{problem}
数列 \(\{a_n\}\)，\(a_n=n+c\)，\(S_7<0\)，\(c\) 的取值范围为\fillinblank{} .
\end{problem}

\begin{answer}
\(\left(-\infty,-4\right)\)
\end{answer}

\begin{solution}
由 \(a_{n}=n+c\) 得 \(\left\{a_{n}\right\}\) 为等差数列，前 \(7\) 项和为 \(S_{7}=\sum_{n=1}^{7}\left(n+c\right)=28+7c\)，亦可用等差中项写成 \(S_{7}=7a_{4}=7\left(4+c\right)\)，由 \(S_{7}<0\) 得 \(4+c<0\)，即 \(c<-4\)，故取值范围为 \(\left(-\infty,-4\right)\)，边界 \(c=-4\) 时 \(S_{7}=0\) 不合要求.
\end{solution}

\begin{problem}
三角形三边长为 \(5\)，\(6\)，\(7\)，则以边长为 \(6\) 的两个顶点为焦点，过另外一个顶点的双曲线的离心率为\fillinblank{} .
\end{problem}

\begin{answer}
\(3\)
\end{answer}

\begin{solution}
以边长为 \(6\) 的两个顶点为焦点，则焦距 \(2c=6\)，得 \(c=3\)，过另一顶点的一支上该点到两焦点的距离分别为 \(5\)，\(7\)，由双曲线定义实轴长 \(2a=\left|7-5\right|=2\)，得 \(a=1\)，故离心率 \(e=\frac{c}{a}=3\).
\end{solution}

\begin{problem}
已知 \(f(x)=x^2\)，\(g(x)=\displaystyle\begin{cases}
f(x), & x\ge 0,\\
-f(-x), & x<0,
\end{cases}\) 求 \(g(x)\le 2-x\) 的 \(x\) 的取值范围\fillinblank{} .
\end{problem}

\begin{answer}
\(\left(-\infty,1\right]\)
\end{answer}

\begin{solution}
由 \(f\left(x\right)=x^{2}\) 得 \(f\left(-x\right)=\left(-x\right)^{2}=x^{2}\)，故 \(x\ge 0\) 时 \(g\left(x\right)=x^{2}\)，\(x<0\) 时 \(g\left(x\right)=-x^{2}\)，当 \(x\ge 0\) 时 \(x^{2}\le 2-x\) 化为 \(x^{2}+x-2\le 0\)，即 \(\left(x+2\right)\left(x-1\right)\le 0\)，交 \(x\ge 0\) 得 \(0\le x\le 1\)，当 \(x<0\) 时 \(-x^{2}\le 2-x\) 化为 \(x^{2}-x+2\ge 0\)，判别式 \(\Delta=\left(-1\right)^{2}-8=-7<0\) 且二次项系数为正，故对一切 \(x<0\) 恒成立，合并得 \(x\) 的取值范围为 \(\left(-\infty,1\right]\)，经端点 \(x=1\) 回代取等号成立.
\end{solution}

\begin{problem}
在棱柱 \(ABCD-A_1B_1C_1D_1\) 中，底面 \(ABCD\) 为平行四边形，\(\left|AA_1\right|=3\)，\(\left|BD\right|=4\)，\(\overrightarrow{AB_1}\cdot\overrightarrow{BC}-\overrightarrow{AD_1}\cdot\overrightarrow{DC}=5\)，设异面直线 \(AA_1\) 与 \(BD\) 的夹角为 \(\theta\)，则 \(\cos\theta=\)\fillinblank{} .
\end{problem}

\begin{answer}
\(\frac{5}{12}\)
\end{answer}

\begin{solution}
记 \(\overrightarrow{AB}=\bs u\)，\(\overrightarrow{AD}=\bs v\)，\(\overrightarrow{AA_{1}}=\bs w\)，由棱柱性质 \(\overrightarrow{BC}=\bs v\)，\(\overrightarrow{DC}=\bs u\)，\(\overrightarrow{AB_{1}}=\bs u+\bs w\)，\(\overrightarrow{AD_{1}}=\bs v+\bs w\)，于是 \(\overrightarrow{AB_{1}}\cdot\overrightarrow{BC}-\overrightarrow{AD_{1}}\cdot\overrightarrow{DC}=\left(\bs u+\bs w\right)\cdot\bs v-\left(\bs v+\bs w\right)\cdot\bs u=\bs w\cdot\left(\bs v-\bs u\right)=5\)，又 \(\overrightarrow{BD}=\bs v-\bs u\)，\(\left|BD\right|=\left|\bs v-\bs u\right|=4\)，\(\left|\bs w\right|=\left|AA_{1}\right|=3\)，异面直线夹角 \(\theta\in\left[0,\frac{\pi}{2}\right]\) 满足 \(\cos\theta=\frac{\left|\bs w\cdot\left(\bs v-\bs u\right)\right|}{\left|\bs w\right|\left|\bs v-\bs u\right|}=\frac{5}{3\cdot 4}=\frac{5}{12}\).
\end{solution}

\begin{problem}
正方形草地 \(ABCD\) 边长 \(1.2\)，\(E\) 到 \(AB\)，\(AD\) 距离为 \(0.2\)，\(F\) 到 \(BC\)，\(CD\) 距离为 \(0.4\)，有个位于正方形草地 \(ABCD\) 内的圆形通道经过 \(E\)，\(F\)，且与 \(AD\) 相切于 \(AD\) 边上一点，求圆形通道的周长\fillinblank{} .（精确到 \(0.01\)）

\bitmapfigure[max width=0.28\linewidth]{img/2024/shanghai_spring/q11_fig1.png}
\end{problem}

\begin{answer}
\(2.73\)
\end{answer}

\begin{solution}
以 \(A(0,0)\)，\(B(1.2,0)\)，\(C(1.2,1.2)\)，\(D(0,1.2)\) 建立坐标系，则 \(E(0.2,0.2)\)，\(F(0.8,0.8)\). \(E\)，\(F\) 连线的垂直平分线为 \(x+y=1\)，故圆心必为 \(O(t,1-t)\)，半径为 \(r(t)=\sqrt{\left(t-0.2\right)^{2}+\left(0.8-t\right)^{2}}\). 由修订题干圆与 \(AD\) 相切于 \(AD\) 边上一点，\(AD\) 为 \(x=0\)，故 \(r(t)=\left|t\right|\)，即 \(t^{2}=\left(t-0.2\right)^{2}+\left(0.8-t\right)^{2}\)，化为 \(t^{2}-2t+0.68=0\)，得 \(t=1\pm\sqrt{0.32}\). 其中大根 \(t=1+\sqrt{0.32}\approx 1.57\) 给出圆心在正方形外，舍去；小根 \(t=1-\sqrt{0.32}\approx 0.43\) 给出圆心在正方形内，切点 \(\left(0,1-t\right)\approx\left(0,0.57\right)\) 落在 \(AD\) 边上，且圆到 \(AB\)，\(BC\)，\(CD\) 的距离均不小于半径，满足位于草地内的要求，故取 \(r=t\approx 0.43\). 于是周长为 \(2\pi r\approx 2.73\).
\end{solution}

\begin{problem}
\(a_1=2\)，\(a_2=4\)，\(a_3=8\)，\(a_4=16\)，任意 \(b_1\)，\(b_2\)，\(b_3\)，\(b_4\in \R\)，满足 \(\{a_i+a_j\mid 1\le i<j\le 4\}=\{b_i+b_j\mid 1\le i<j\le 4\}\)，求有序数列 \(\{b_1,b_2,b_3,b_4\}\) 有\fillinblank{} 对.
\end{problem}

\begin{answer}
\(48\)
\end{answer}

\begin{solution}
记\(a\)的六个两两和为\(6\)，\(10\)，\(12\)，\(18\)，\(20\)，\(24\)，共六个不同值，故\(\{b_{i}+b_{j}\}\)作为集合有六元，从而四个\(b_{i}\)两两不同，否则必有两个和重合. 将\(b_{i}\)按从小到大记为\(c_{1}<c_{2}<c_{3}<c_{4}\)，则最小和只能是\(c_{1}+c_{2}\)，最大和只能是\(c_{3}+c_{4}\)，故\(c_{1}+c_{2}=6\)，\(c_{3}+c_{4}=24\). 六个和之和为\(3(c_{1}+c_{2}+c_{3}+c_{4})=90\)，故\(c_{1}+c_{2}+c_{3}+c_{4}=30\). 在剩下的四个和中，最小者为\(c_{1}+c_{3}\)，最大者为\(c_{2}+c_{4}\)，故\(c_{1}+c_{3}=10\)，\(c_{2}+c_{4}=20\)，剩下\(\{c_{1}+c_{4},c_{2}+c_{3}\}=\{12,18\}\). 若\(c_{1}+c_{4}=12\)，\(c_{2}+c_{3}=18\)，解得\(c_{1}=-1\)，\(c_{2}=7\)，\(c_{3}=11\)，\(c_{4}=13\). 若\(c_{1}+c_{4}=18\)，\(c_{2}+c_{3}=12\)，解得\(c_{1}=2\)，\(c_{2}=4\)，\(c_{3}=8\)，\(c_{4}=16\). 直接验算两组的有序对和集均为\(\{6,10,12,18,20,24\}\). 故无序集合恰两种，每种有\(4!=24\)个有序排列，有序数列共\(48\)对.
\end{solution}
\section{单选题}

\begin{problem}
已知 \(a\)，\(b\)，\(c\in \R\)，\(b>c\)，下列不等式恒成立的是（\quad）.

\choices
    {\(a+b^2>a+c^2\)}
    {\(a^2-c>a^2-b\)}
    {\(ab^2>ac^2\)}
    {\(a^2b>a^2c\)}
\end{problem}

\begin{answer}
B
\end{answer}

\begin{solution}
因为\(b>c\)，所以\(-c>-b\)，故\(a^{2}-c>a^{2}-b\)对一切\(a\)，\(b\)，\(c\in\R\)恒成立，B正确. \(b>c\)不能推出\(b^{2}>c^{2}\)，例如\(b=0\)，\(c=-1\)时\(b^{2}<c^{2}\). 取\(a=0\)，\(b=0\)，\(c=-1\)，则\(a+b^{2}=0\)，\(a+c^{2}=1\)，A不成立；\(ab^{2}=0\)，\(ac^{2}=0\)，C不成立；\(a^{2}b=0\)，\(a^{2}c=0\)，D不成立. 故排除A，C，D.
\end{solution}

\begin{problem}
空间中有两个不同的平面 \(\alpha\)，\(\beta\) 和两条不同的直线 \(m\)，\(n\)，则下列说法中正确的是（\quad）.

\choices
    {若 \(\alpha\perp\beta\)，\(m\perp\alpha\)，\(m\perp n\)，则 \(n\perp\beta\)}
    {若 \(\alpha\perp\beta\)，\(m\perp\alpha\)，\(n\perp\beta\)，则 \(m\perp n\)}
    {若 \(\alpha\parallel\beta\)，\(m\parallel\alpha\)，\(n\parallel\beta\)，则 \(m\parallel n\)}
    {若 \(\alpha\parallel\beta\)，\(m\parallel\alpha\)，\(m\parallel n\)，则 \(n\parallel\beta\)}
\end{problem}

\begin{answer}
B
\end{answer}

\begin{solution}
设\(\bs{u}\)，\(\bs{v}\)分别为\(\alpha\)，\(\beta\)的法向量，则\(\alpha\perp\beta\)等价于\(\bs{u}\perp\bs{v}\). 若\(m\perp\alpha\)，\(n\perp\beta\)，则\(m\)的方向向量平行于\(\bs{u}\)，\(n\)的方向向量平行于\(\bs{v}\)，故\(m\perp n\)，B正确. 取\(\alpha\)为平面\(x=0\)，\(\beta\)为平面\(y=0\)，\(m\)为\(x\)轴，则\(m\perp\alpha\). 过原点沿方向\((0,1,1)\)作\(n\)，则\(m\perp n\)且\(n\subset\alpha\)，但\(n\)不垂直于\(\beta\)，故A错误. 取\(\alpha\)为平面\(z=0\)，\(\beta\)为平面\(z=1\)，\(m\)沿方向\((1,0,0)\)，\(n\)沿方向\((0,1,0)\)，则\(m\parallel\alpha\)，\(n\parallel\beta\)，但\(m\)与\(n\)不平行，故C错误. 取同样的平行平面，\(m\)为直线\(y=0\)，\(z=2\)，\(n\)为直线\(y=0\)，\(z=1\)，则\(m\parallel\alpha\)，\(m\parallel n\)，\(\alpha\parallel\beta\)，但\(n\subset\beta\)，按直线在平面外才算平行计\(n\parallel\beta\)不成立，故D错误.
\end{solution}

\begin{problem}
有四种礼盒，前三种里面分别仅装有中国结，记事本，笔袋，第四个礼盒里面三种礼品都有，现从中任选一个盒子，设事件 \(A\)：所选盒中有中国结，事件 \(B\)：所选盒中有记事本，事件 \(C\)：所选盒中有笔袋，则（\quad）.

\choices
    {事件 \(A\) 与事件 \(B\) 互斥}
    {事件 \(A\) 与事件 \(B\) 相互独立}
    {事件 \(A\) 与事件 \(B\cup C\) 互斥}
    {事件 \(A\) 与事件 \(B\cap C\) 相互独立}
\end{problem}

\begin{answer}
B
\end{answer}

\begin{solution}
从四个盒子中任选一个是等可能的. 含中国结的是第一个盒与第四个盒，故\(P(A)=\frac{2}{4}=\frac{1}{2}\). 含记事本的是第二个盒与第四个盒，故\(P(B)=\frac{1}{2}\). 同时含两者的是第四个盒，故\(P(A\cap B)=\frac{1}{4}=P(A)P(B)\)，且\(A\cap B\ne\varnothing\)，所以\(A\)与\(B\)不互斥但相互独立. 含记事本或笔袋的是第二、三、四个盒，故\(P(B\cup C)=\frac{3}{4}\)，而\(P(A\cap(B\cup C))=\frac{1}{4}\ne P(A)P(B\cup C)\)，且交集非空，故选项 C 不成立. 同时含记事本与笔袋的只有第四个盒，故\(P(B\cap C)=\frac{1}{4}\)，\(P(A\cap(B\cap C))=\frac{1}{4}\ne P(A)P(B\cap C)\)，故选项 D 不成立.
\end{solution}

\begin{problem}
现定义如下：当 \(x\in(n,n+1)\) 时（\(n\in \N\)），若 \(f(x+1)=f'(x)\)，则称 \(f(x)\) 为延展函数. 已知当 \(x\in(0,1)\) 时，\(g(x)=\e^x\)，且 \(h(x)=x^{10}\)，且 \(g(x)\)，\(h(x)\) 均为延展函数，则以下结论（\quad）.

\begin{circlelist}
\item 存在 \(y=kx+b\)（\(k\)，\(b\in\R\)，\(k\)，\(b\ne 0\)）与 \(y=g(x)\) 有无穷多个交点；
\item 存在 \(y=kx+b\)（\(k\)，\(b\in\R\)，\(k\)，\(b\ne 0\)）与 \(y=h(x)\) 有无穷多个交点.
\end{circlelist}

\choices
    {（1）（2）都成立}
    {（1）（2）都不成立}
    {（1）成立（2）不成立}
    {（1）不成立（2）成立}
\end{problem}

\begin{answer}
D
\end{answer}

\begin{solution}
当\(x\in(0,1)\)时\(g(x)=\e^{x}\). 对\(x\in(n,n+1)\)，\(n\ge 1\)，有\(g(x)=g^{\prime}(x-1)\). 归纳得\(x\in(n,n+1)\)时\(g(x)=\e^{x-n}\)，值域恒在\((1,\e)\)内. 记\(\varphi(x)=\e^{x-n}-kx-b\)，则\(\varphi^{\prime\prime}(x)>0\)，故每段至多两个交点. 若\(k>0\)，当\(x>(\e-b)/k\)时\(kx+b>\e\ge g(x)\). 若\(k<0\)，当\(x\)充分大时\(kx+b<1\le g(x)\). 于是\(k\ne 0\)的直线只可能在有限多段内与\(g\)相交，交点有限，（1）不成立. 当\(x\in(0,1)\)时\(h(x)=x^{10}\)，归纳得\(0\le n\le 10\)且\(x\in(n,n+1)\)时\(h(x)=\frac{10!}{(10-n)!}(x-n)^{10-n}\). 特别在\((9,10)\)上\(h(x)=3628800(x-9)\). 取\(y=3628800x-32659200\)，则\(k\ne 0\)，\(b\ne 0\)，且该直线在整个\((9,10)\)上与\(h\)重合，有无穷多个交点，故（2）成立. 于是（1）不成立（2）成立.
\end{solution}
\section{解答题}

\begin{problem}
已知 \(f(x)=\sin\left(\omega x+\frac{\pi}{3}\right)\)，\(\omega>0\).

\begin{enumerate}
\item 设 \(\omega=1\)，求 \(y=f(x)\)，\(x\in[0,\pi]\) 的值域；
\item \(a>\pi\)（\(a\in\R\)），\(f(x)\) 的最小正周期为 \(\pi\)，若在 \(x\in[\pi,a]\) 上恰有 \(3\) 个零点，求 \(a\) 的取值范围.
\end{enumerate}
\end{problem}

\begin{solution}
\begin{enumerate}
\item 当 \(\omega=1\) 时，\(f(x)=\sin\left(x+\frac{\pi}{3}\right)\). 由 \(x\in\left[0,\pi\right]\) 得 \(t=x+\frac{\pi}{3}\in\left[\frac{\pi}{3},\frac{4\pi}{3}\right]\)，\(y=\sin t\) 在 \(\left[\frac{\pi}{3},\frac{\pi}{2}\right]\) 上单调递增，在 \(\left[\frac{\pi}{2},\frac{4\pi}{3}\right]\) 上单调递减. 因为 \(\sin\frac{\pi}{2}=1\)，\(\sin\frac{\pi}{3}=\frac{\sqrt3}{2}\)，\(\sin\frac{4\pi}{3}=-\frac{\sqrt3}{2}\)，所以值域为 \(\left[-\frac{\sqrt3}{2},1\right]\).
\item 由最小正周期 \(\frac{2\pi}{\omega}=\pi\) 得 \(\omega=2\)，\(f(x)=\sin\left(2x+\frac{\pi}{3}\right)\). 令 \(f(x)=0\) 得 \(2x+\frac{\pi}{3}=k\pi\)（\(k\in\Z\)），即 \(x=-\frac{\pi}{6}+\frac{k\pi}{2}\). 在 \(x\ge\pi\) 的零点中，\(k=2\) 时 \(x=\frac{5\pi}{6}<\pi\) 应舍去，\(k=3\)，\(k=4\)，\(k=5\) 依次得 \(x=\frac{4\pi}{3}\)，\(x=\frac{11\pi}{6}\)，\(x=\frac{7\pi}{3}\)，下一个零点 \(k=6\) 时 \(x=\frac{17\pi}{6}\). 故在 \(\left[\pi,a\right]\) 上恰有 \(3\) 个零点当且仅当 \(a\in\left[\frac{7\pi}{3},\frac{17\pi}{6}\right)\).
\end{enumerate}
\end{solution}

\begin{problem}
如图，\(PA\)，\(PB\)，\(PC\) 为圆锥三条母线，\(AB=AC\).

\begin{enumerate}
\item 证明：\(PA\perp BC\)；
\item 若圆锥侧面积为 \(\sqrt3\pi\)，\(BC\) 为底面直径，\(BC=2\)，求二面角 \(B-PA-C\) 的大小.
\end{enumerate}

\FigureLayoutDeclare{img/2024/shanghai_spring/q18_fig1.png}{8.5cm}{0bp 9.002bp 0bp 3.001bp}
\bitmapfigure[max width=0.34\linewidth]{img/2024/shanghai_spring/q18_fig1.png}
\end{problem}

\begin{solution}
\begin{enumerate}
\item 取 \(BC\) 的中点 \(O\)，连接 \(AO\)，\(PO\). 因为 \(AB=AC\)，所以 \(AO\perp BC\). 因为 \(PA\)，\(PB\)，\(PC\) 均为圆锥母线，所以 \(PB=PC\)，故 \(PO\perp BC\). 又 \(AO\cap PO=O\)，\(AO\subset\) 平面 \(PAO\)，\(PO\subset\) 平面 \(PAO\)，所以 \(BC\perp\) 平面 \(PAO\). 因为 \(PA\subset\) 平面 \(PAO\)，所以 \(PA\perp BC\).
\item 因为 \(BC\) 为底面直径，所以 \(O\) 为底面圆心，\(PO\perp\) 底面 \(BAC\). 又 \(AB=AC\) 得 \(AO\perp BC\). 以 \(O\) 为原点，分别以 \(\overrightarrow{OB}\) 的方向为 \(x\) 轴正方向，以 \(\overrightarrow{OA}\) 的方向为 \(y\) 轴正方向，以 \(\overrightarrow{OP}\) 的方向为 \(z\) 轴正方向建立空间直角坐标系. 圆锥侧面积 \(\pi rl=\sqrt3\pi\)，底面半径 \(r=1\)，故母线 \(l=\sqrt3\)，即 \(PA=\sqrt3\). 又 \(OA=1\)，所以高 \(PO=\sqrt{PA^2-OA^2}=\sqrt2\). 于是 \(P(0,0,\sqrt2)\)，\(A(0,1,0)\)，\(B(1,0,0)\)，\(C(-1,0,0)\)，\(\overrightarrow{PA}=(0,1,-\sqrt2)\)，\(\overrightarrow{PB}=(1,0,-\sqrt2)\)，\(\overrightarrow{PC}=(-1,0,-\sqrt2)\). 设平面 \(PAB\) 的法向量 \(\overrightarrow{n}_1=(x_1,y_1,z_1)\)，则 \(\overrightarrow{n}_1\cdot\overrightarrow{PA}=y_1-\sqrt2z_1=0\)，\(\overrightarrow{n}_1\cdot\overrightarrow{PB}=x_1-\sqrt2z_1=0\)，取 \(z_1=1\) 得 \(\overrightarrow{n}_1=(\sqrt2,\sqrt2,1)\). 设平面 \(PAC\) 的法向量 \(\overrightarrow{n}_2=(x_2,y_2,z_2)\)，则 \(\overrightarrow{n}_2\cdot\overrightarrow{PA}=y_2-\sqrt2z_2=0\)，\(\overrightarrow{n}_2\cdot\overrightarrow{PC}=-x_2-\sqrt2z_2=0\)，取 \(z_2=1\) 得 \(\overrightarrow{n}_2=(-\sqrt2,\sqrt2,1)\). 于是 \(\cos\left\langle\overrightarrow{n}_1,\overrightarrow{n}_2\right\rangle=\frac{1}{\sqrt5\cdot\sqrt5}=\frac15\). 取 \(PA\) 的中点 \(M\left(0,\frac12,\frac{\sqrt2}{2}\right)\)，向量 \(\overrightarrow{u}_B=\left(1,-\frac23,-\frac{\sqrt2}{3}\right)\) 在平面 \(PAB\) 内且垂直于 \(PA\)，向量 \(\overrightarrow{u}_C=\left(-1,-\frac23,-\frac{\sqrt2}{3}\right)\) 在平面 \(PAC\) 内且垂直于 \(PA\)，它们的夹角余弦为 \(-\frac15\)，故二面角 \(B-PA-C\) 为钝角，其大小为 \(\pi-\arccos\frac15\).
\end{enumerate}
\end{solution}

\begin{problem}
水果分为一级果和二级果，共 \(136\) 箱，其中一级果 \(102\) 箱，二级果 \(34\) 箱.

\begin{enumerate}
\item 随机挑选两箱水果，求恰好一级果和二级果各一箱的概率；
\item 进行分层抽样，共抽 \(8\) 箱水果，求一级果和二级果各几箱；
\item 抽取若干箱水果，其中一级果共 \(120\) 个，单果质量平均数为 \(303.45\) 克，方差为 \(603.46\)；二级果 \(48\) 个，单果质量平均数为 \(240.41\) 克，方差为 \(648.21\)；求 \(168\) 个水果的方差和平均数，并预估果园中单果的质量.
\end{enumerate}
\end{problem}

\begin{solution}
\begin{enumerate}
\item 设 \(A\) 为恰好取到一级果和二级果各一箱，则 \(P(A)=\frac{\mathrm C_{102}^1\mathrm C_{34}^1}{\mathrm C_{136}^2}=\frac{3468}{9180}=\frac{17}{45}\).
\item 一级果与二级果箱数比为 \(102:34=3:1\)，故 \(8\) 箱中一级果抽 \(8\times\frac{3}{3+1}=6\) 箱，二级果抽 \(8\times\frac{1}{3+1}=2\) 箱.
\item 记一级果单果质量均值与方差为 \(\bar x=303.45\text{ 克}\)，\(S_x^2=603.46\text{ 克}^2\)，二级果单果质量均值与方差为 \(\bar y=240.41\text{ 克}\)，\(S_y^2=648.21\text{ 克}^2\). 总体均值 \(\bar z=\frac{120}{168}\bar x+\frac{48}{168}\bar y=285.44\text{ 克}\)，总体方差 \(S^2=\frac{120}{168}\left[S_x^2+(\bar x-\bar z)^2\right]+\frac{48}{168}\left[S_y^2+(\bar y-\bar z)^2\right]=1427.27\text{ 克}^2\). 按果园中两类果箱数占比 \(\frac{102}{136}\)，\(\frac{34}{136}\) 预估单果质量为 \(\frac{102}{136}\bar x+\frac{34}{136}\bar y=287.69\text{ 克}\).
\end{enumerate}
\end{solution}

\begin{problem}
在平面直角坐标系 \(xOy\) 中，已知点 \(A\) 为椭圆 \(\Gamma:\frac{x^2}{6}+\frac{y^2}{2}=1\) 上一点，\(F_1\)，\(F_2\) 分别为椭圆的左，右焦点.

\begin{enumerate}
\item 若点 \(A\) 的横坐标为 \(2\)，求 \(\left|AF_1\right|\) 的长；
\item 设 \(\Gamma\) 的上，下顶点分别为 \(M_1\)，\(M_2\)，记 \(\triangle AF_1F_2\) 的面积为 \(S_1\)，\(\triangle AM_1M_2\) 的面积为 \(S_2\)，若 \(S_1\ge S_2\)，求 \(\left|OA\right|\) 的取值范围；
\item 若点 \(A\) 在 \(x\) 轴上方，设直线 \(AF_2\) 与 \(\Gamma\) 交于点 \(B\)，与 \(y\) 轴交于点 \(K\)，\(KF_1\) 延长线与 \(\Gamma\) 交于点 \(C\)，是否存在 \(x\) 轴上方的点 \(C\)，使得 \(\overrightarrow{F_1A}+\overrightarrow{F_1B}+\overrightarrow{F_1C} =\lambda\left(\overrightarrow{F_2A}+\overrightarrow{F_2B}+\overrightarrow{F_2C}\right)\)（\(\lambda\in\R\)）
成立？若存在，请求出点 \(C\) 的坐标；若不存在，请说明理由.
\end{enumerate}
\end{problem}

\begin{solution}
\begin{enumerate}
\item 由 \(a^2=6\)，\(b^2=2\) 得 \(c=\sqrt{a^2-b^2}=2\)，故 \(F_1\left(-2,0\right)\)，\(F_2\left(2,0\right)\)，离心率 \(e=\frac{c}{a}=\frac{\sqrt{6}}{3}\)，长轴长 \(2a=2\sqrt{6}\). 设 \(A\left(2,y\right)\)，代入椭圆得 \(\frac{4}{6}+\frac{y^2}{2}=1\)，即 \(y^2=\frac{2}{3}\). 由两点距离得 \(\left|AF_1\right|=\sqrt{\left(2+2\right)^2+y^2}=\sqrt{16+\frac{2}{3}}=\frac{5\sqrt{6}}{3}\). 用焦半径独立验算，左焦点半径 \(\left|AF_1\right|=a+ex_1=\sqrt{6}+\frac{\sqrt{6}}{3}\times 2=\frac{5\sqrt{6}}{3}\)，与上式一致，由定义 \(\left|AF_1\right|+\left|AF_2\right|=2a\) 亦得同一结果，所以 \(\left|AF_1\right|=\frac{5\sqrt{6}}{3}\).
\item 记 \(A\left(x,y\right)\)，\(M_1\left(0,\sqrt{2}\right)\)，\(M_2\left(0,-\sqrt{2}\right)\)，则 \(\left|F_1F_2\right|=4\)，\(\left|M_1M_2\right|=2\sqrt{2}\). 若 \(x=0\) 则 \(A\)，\(M_1\)，\(M_2\) 共线不成三角形，若 \(y=0\) 则 \(A\)，\(F_1\)，\(F_2\) 共线不成三角形，故三角形面积要求 \(xy\ne 0\). 此时 \(S_1=\frac{1}{2}\left|F_1F_2\right|\cdot \left|y\right|=2\left|y\right|\)，\(S_2=\frac{1}{2}\left|M_1M_2\right|\cdot \left|x\right|=\sqrt{2}\left|x\right|\). 由 \(S_1\ge S_2\) 得 \(2\left|y\right|\ge \sqrt{2}\left|x\right|\)，即 \(2y^2\ge x^2\). 结合 \(\frac{x^2}{6}+\frac{y^2}{2}=1\) 得 \(x^2=6-3y^2\)，故 \(2y^2\ge 6-3y^2\)，即 \(y^2\ge \frac{6}{5}\). 又椭圆上 \(y^2\le 2\)，而 \(x\ne 0\) 得 \(y^2<2\)，所以 \(\frac{6}{5}\le y^2<2\). 而 \(\left|OA\right|^2=x^2+y^2=6-2y^2\)，故 \(2<\left|OA\right|^2\le \frac{18}{5}\)，即 \(\left|OA\right|\in \left(\sqrt{2},\frac{3\sqrt{10}}{5}\right]\). 右端在 \(y^2=\frac{6}{5}\) 取到，左端 \(x=0\) 按面积公式形式上满足不等式但不成三角形故舍去，与原答案左开右闭一致.
\item 设 \(A\left(x_1,y_1\right)\)，\(y_1>0\)，\(B\left(x_2,y_2\right)\)，\(C\left(x_C,y_C\right)\)，\(y_C>0\)，\(F_1\left(-2,0\right)\)，\(F_2\left(2,0\right)\). 记 \(\overrightarrow{F_1A}+\overrightarrow{F_1B}+\overrightarrow{F_1C}=\left(x_1+x_2+x_C+6,y_1+y_2+y_C\right)\)，\(\overrightarrow{F_2A}+\overrightarrow{F_2B}+\overrightarrow{F_2C}=\left(x_1+x_2+x_C-6,y_1+y_2+y_C\right)\). 两向量共线等价于 \(\left(x_1+x_2+x_C+6\right)\left(y_1+y_2+y_C\right)=\left(x_1+x_2+x_C-6\right)\left(y_1+y_2+y_C\right)\)，即 \(12\left(y_1+y_2+y_C\right)=0\)，故必 \(y_1+y_2+y_C=0\)，且此时两向量均为水平向量，只要 \(\overrightarrow{F_2A}+\overrightarrow{F_2B}+\overrightarrow{F_2C}\ne \bs{0}\) 即有 \(\lambda\) 成立. 若 \(x_1=2\) 则直线 \(AF_2\) 为 \(x=2\) 无 \(y\) 轴交点，故舍去，设 \(AF_2\) 为 \(x=my+2\)，其中 \(m=\frac{x_1-2}{y_1}\ne 0\)，则 \(K\left(0,-\frac{2}{m}\right)\). 过 \(K\)，\(F_1\) 的直线为 \(x=-my-2\). 将两直线分别代入 \(x^2+3y^2=6\) 均得 \(\left(m^2+3\right)y^2+4my-2=0\)，其两根之积为 \(\frac{-2}{m^2+3}<0\)，一正一负. 因为椭圆关于 \(y\) 轴对称，直线 \(x=-my-2\) 与椭圆的两交点恰为直线 \(x=my+2\) 与椭圆两交点关于 \(y\) 轴的镜像，即 \(\left(-x_1,y_1\right)\)，\(\left(-x_2,y_2\right)\). 其中纵坐标为正者为上交点，故 \(C\left(-x_1,y_1\right)\) 得证. 于是 \(y_1+y_2+y_C=2y_1+y_2=0\)，即 \(y_2+2y_1=0\). 由韦达 \(y_1+y_2=-\frac{4m}{m^2+3}\)，\(y_1y_2=\frac{-2}{m^2+3}\)，代入 \(y_2=-2y_1\) 得 \(-y_1=-\frac{4m}{m^2+3}\)，\(-2y_1^2=\frac{-2}{m^2+3}\)，即 \(y_1^2=\frac{1}{m^2+3}\)，\(y_1=\frac{4m}{m^2+3}\). 由 \(y_1>0\) 得 \(m>0\)，解得 \(m^2=\frac{1}{5}\)，\(m=\frac{\sqrt{5}}{5}\)，\(y_1=\frac{\sqrt{5}}{4}\)，\(x_1=my_1+2=\frac{9}{4}\). 故 \(C\left(-\frac{9}{4},\frac{\sqrt{5}}{4}\right)\). 回代得 \(B\left(\frac{3}{2},-\frac{\sqrt{5}}{2}\right)\) 在椭圆上，\(K\left(0,-2\sqrt{5}\right)\)，\(\overrightarrow{F_2A}+\overrightarrow{F_2B}+\overrightarrow{F_2C}=\left(-\frac{9}{2},0\right)\ne \bs{0}\)，\(\overrightarrow{F_1A}+\overrightarrow{F_1B}+\overrightarrow{F_1C}=\left(\frac{15}{2},0\right)\)，取 \(\lambda=-\frac{5}{3}\) 成立. 所以存在 \(x\) 轴上方的点 \(C\left(-\frac{9}{4},\frac{\sqrt{5}}{4}\right)\) 满足要求.
\end{enumerate}
\end{solution}

\begin{problem}
记 \(M(a)=\{t\mid t=f(x)-f(a),x\ge a\}\)，\(L(a)=\{t\mid t=f(x)-f(a),x\le a\}\).

\begin{enumerate}
\item 若 \(f(x)=x^2+1\)，求 \(M(1)\) 和 \(L(1)\)；
\item 若 \(f(x)=x^3-3x^2\)，求证：对于任意 \(a\in\R\)，都有 \(M(a)\subseteq[-4,+\infty)\)，且存在 \(a\)，使得 \(-4\in M(a)\).
\item 已知定义在 \(\R\) 上的 \(f(x)\) 有最小值，求证“\(f(x)\) 是偶函数”的充要条件是“对于任意正实数 \(c\)，均有 \(M(-c)=L(c)\)”.
\end{enumerate}
\end{problem}

\begin{solution}
\begin{enumerate}
\item 因为 \(f(x)=x^2+1\)，所以 \(f(1)=2\)，\(M(1)=\{x^2-1\mid x\ge 1\}\)，\(L(1)=\{x^2-1\mid x\le 1\}\).当 \(x\ge 1\) 时，\(x^2-1\) 严格递增并趋于 \(+\infty\)，最小值为 \(0\)，故 \(M(1)=\left[0,+\infty\right)\).当 \(x\le 1\) 时，对一切 \(y\ge 0\) 都有 \(-\sqrt{y}\le 1\)，故 \(x^2\) 取遍 \(\left[0,+\infty\right)\)，从而 \(L(1)=\left[-1,+\infty\right)\).

\item 记 \(g(x)=x^3-3x^2-a^3+3a^2\)，\(x\ge a\)，则 \(M(a)\) 为 \(g\) 在 \(\left[a,+\infty\right)\) 上的值域，且 \(g(a)=0\)，\(\displaystyle\lim_{x\to +\infty}g(x)=+\infty\).由 \(g'(x)=3x^2-6x=3x\left(x-2\right)\) 得驻点 \(0\)，\(2\)；\(g\) 在 \(\left(-\infty,0\right)\)，\(\left(2,+\infty\right)\) 上严格递增，在 \(\left(0,2\right)\) 上严格递减.

当 \(a\ge 2\) 时，\(g\) 在 \(\left[a,+\infty\right)\) 上严格递增，故最小值为 \(g(a)=0\)，\(M(a)=\left[0,+\infty\right)\subseteq\left[-4,+\infty\right)\).

当 \(0\le a<2\) 时，\(g\) 在 \(\left[a,2\right]\) 上严格递减，在 \(\left[2,+\infty\right)\) 上严格递增，故最小值在 \(x=2\) 处取得，\(t_{\min}=g(2)=-a^3+3a^2-4\).因为 \(-a^3+3a^2=a^2\left(3-a\right)\ge 0\)，等号仅在 \(a=0\) 处成立，所以 \(t_{\min}\ge -4\)，\(M(a)=\left[t_{\min},+\infty\right)\subseteq\left[-4,+\infty\right)\).

当 \(a<0\) 时，\(g\) 在 \(\left[a,0\right]\) 上严格递增，在 \(\left[0,2\right]\) 上严格递减，在 \(\left[2,+\infty\right)\) 上严格递增.此时 \(g(0)=-a^3+3a^2=a^2\left(3-a\right)>0\)，\(g\) 在 \(\left[a,0\right]\) 上的值域为 \(\left[0,g(0)\right]\)，在 \(\left[0,+\infty\right)\) 上的值域为 \(\left[g(2),+\infty\right)\)，其中 \(g(2)=g(0)-4\).记 \(h(a)=g(2)=-a^3+3a^2-4\)，则 \(h'(a)=-3a^2+6a=3a\left(2-a\right)<0\)，故 \(h\) 在 \(\left(-\infty,0\right)\) 上严格递减，从而 \(h(a)>h(0)=-4\).于是 \(M(a)=\left[t_{\min},+\infty\right)\)，其中 \(t_{\min}=\min\{0,g(2)\}>-4\)，故 \(M(a)\subseteq\left[-4,+\infty\right)\).

综上，对任意 \(a\in\R\) 都有 \(M(a)\subseteq\left[-4,+\infty\right)\).取 \(a=0\)，则 \(f(0)=0\)，\(f(2)=-4\) 且 \(2\ge 0\)，故 \(-4=f(2)-f(0)\in M(0)\)，存在性成立.

\item 先证必要性.设 \(f\) 为偶函数，\(c>0\)，则 \(f(-c)=f(c)\)，\(M(-c)=\{f(x)-f(-c)\mid x\ge -c\}\)，\(L(c)=\{f(x)-f(c)\mid x\le c\}\).对任意 \(t\in M(-c)\)，存在 \(x\ge -c\) 使 \(t=f(x)-f(-c)\)，令 \(y=-x\le c\)，则 \(t=f(-y)-f(-c)=f(y)-f(c)\in L(c)\)；反之亦然，故 \(M(-c)=L(c)\).

再证充分性.设 \(f\) 在 \(\R\) 上有全局最小值 \(m\)，存在 \(a\) 使 \(f(a)=m\)，且对一切 \(c>0\) 有 \(M(-c)=L(c)\).当 \(c\ge\left|a\right|\) 时，\(a\in\left[-c,c\right]\)，故 \(M(-c)\) 在 \(x=a\) 处取得最小元 \(m-f(-c)\)，\(L(c)\) 在 \(x=a\) 处取得最小元 \(m-f(c)\).由集合相等得最小元相等，故 \(f(c)=f(-c)\) 对一切 \(c\ge\left|a\right|\) 成立.特别取 \(c=\left|a\right|\) 得 \(f\left(\left|a\right|\right)=f\left(-\left|a\right|\right)\)，而 \(a\) 等于其中之一且函数值为 \(m\)，故 \(\left|a\right|\) 与 \(-\left|a\right|\) 均为全局最小值点.

任取 \(c>0\)，因为 \(\left|a\right|\ge -c\)，所以 \(M(-c)\) 在 \(x=\left|a\right|\) 处取得最小元 \(m-f(-c)\)；因为 \(-\left|a\right|\le c\)，所以 \(L(c)\) 在 \(x=-\left|a\right|\) 处取得最小元 \(m-f(c)\).由 \(M(-c)=L(c)\) 得最小元相等，故 \(f(c)=f(-c)\).于是对一切 \(c\ge 0\) 有 \(f(c)=f(-c)\)，即 \(f\) 为偶函数.
\end{enumerate}
\end{solution}
